\documentclass[a4paper]{article}

% Paquetes básicos
% Español (México) con XeLaTeX: fontspec para fuentes Unicode, polyglossia
% para la división silábica y los nombres del documento en la variante mexicana.
\usepackage{fontspec}
\usepackage{polyglossia}
\setdefaultlanguage[variant=mexican]{spanish}
% Los nombres chinos van como «pinyin (original)»: xeCJK da a los caracteres
% Han su propia fuente; sin ella XeLaTeX los omite con un aviso.
% FandolSong no cubre algunos caracteres de nombres (U+73FA); WenQuanYi Zen Hei sí.
\usepackage[AutoFallBack=true]{xeCJK}
\setCJKmainfont{FandolSong-Regular.otf}
\setCJKfallbackfamilyfont{\CJKrmdefault}{WenQuanYi Zen Hei}
% polyglossia no traduce los nombres de listados.
\AtBeginDocument{\providecommand{\lstlistingname}{}\renewcommand{\lstlistingname}{Listado}%
  \providecommand{\lstlistlistingname}{}\renewcommand{\lstlistlistingname}{Índice de listados}}
\usepackage{amsmath, amssymb}
\usepackage{graphicx}
\usepackage[margin=2.5cm]{geometry}
\usepackage[most]{tcolorbox}
\usepackage{etoolbox}
\usepackage{listings}
\usepackage{hyperref}
\usepackage{booktabs}
\usepackage{subcaption}
\usepackage{float}

% Cajas de resaltado
\newtcolorbox{knowledgebox}[1]{
    enhanced, colback=blue!5!white, colframe=blue!75!black, colbacktitle=blue!75!black,
    coltitle=white, fonttitle=\bfseries, title={#1}, attach boxed title to top left={yshift=-2mm, xshift=2mm},
    boxrule=1pt, sharp corners
}

\newtcolorbox{importantbox}[1]{
    enhanced, colback=yellow!10!white, colframe=yellow!80!black, colbacktitle=yellow!80!black,
    coltitle=black, fonttitle=\bfseries, title={#1}, sharp corners
}

\newtcolorbox{warningbox}[1]{
    enhanced, colback=red!5!white, colframe=red!75!black, colbacktitle=red!75!black,
    coltitle=white, fonttitle=\bfseries, title={#1}, sharp corners
}

% Listados
\lstset{
    language=Python,
    basicstyle=\ttfamily\small,
    keywordstyle=\color{blue},
    stringstyle=\color{red!60!black},
    commentstyle=\color{green!60!black},
    breaklines=true,
    frame=single,
    numbers=left,
    numberstyle=\tiny\color{gray},
    captionpos=b,
    extendedchars=true
}

% Metadatos de la portada
\newcommand{\notetitle}{KAIST CS492D Clase 10: Flow Matching 2}
\newcommand{\noteauthors}{Elaboradas a partir de la clase de Minhyuk Sung}
\newcommand{\notedate}{Otoño de 2025}
\newcommand{\videochannel}{Minhyuk Sung (KAIST)}
\newcommand{\videopublishdate}{2025}
\newcommand{\videoduration}{40 minutos}
\newcommand{\videourl}{https://www.youtube.com/watch?v=o9tccBWdcQ4}
\newcommand{\videocoverpath}{cover.jpg}
\newcommand{\repourl}{https://github.com/hqhq1025/ai-course-notes}

% Pie de página con información del repositorio
\usepackage{fancyhdr}
\pagestyle{fancy}
\fancyhf{}
\fancyfoot[C]{\thepage}
\fancyfoot[R]{\footnotesize\color{gray}\href{\repourl}{AI Course Notes}}
\renewcommand{\headrulewidth}{0pt}
\renewcommand{\footrulewidth}{0pt}

\begin{document}

\begin{titlepage}
\centering
{\Large Notas del curso\par}
\vspace{1.2cm}
{\huge\bfseries \notetitle\par}
\vspace{0.8cm}
{\large \noteauthors\par}
\vspace{0.3cm}
{\large \notedate\par}
\vspace{1.2cm}

\ifdefempty{\videocoverpath}{
{\small Al generar las notas, indique la ruta local de la portada del video.\par}
}{
\includegraphics[width=0.82\textwidth,height=0.45\textheight,keepaspectratio]{\videocoverpath}\par
}

\vfill
\begin{tcolorbox}[width=0.9\textwidth, colback=black!2!white, colframe=black!60, sharp corners]
\textbf{Autor o canal del video}: \videochannel\par
\textbf{Fecha de publicación}: \videopublishdate\par
\textbf{Duración del video}: \videoduration\par
\ifdefempty{\videourl}{}{
\textbf{Enlace del video}: \href{\videourl}{\nolinkurl{\videourl}}\par
}\par
\textbf{Página del proyecto}: \href{\repourl}{\nolinkurl{\repourl}}\par
\end{tcolorbox}
\end{titlepage}

\tableofcontents
\newpage

%% --- Inicio del contenido --- %%

\section{Introducción: de Flow Matching 1 a Flow Matching 2}

Esta clase es la número 10 del curso KAIST CS492D (Modelos de difusión y modelos de flujo) y constituye la segunda parte del módulo sobre Flow Matching. El ponente Minhyuk Sung primero repasa el marco básico de Flow Matching introducido en la clase anterior —los tres conceptos centrales de mapa de flujo (Flow Map), campo vectorial (Vector Field) y trayectoria de probabilidad (Probability Path)—, sus relaciones mutuas y, a partir de esta base, profundiza en los siguientes temas clave:

\begin{enumerate}
    \item La \textbf{equivalencia matemática} entre modelos de difusión y modelos de flujo: bajo qué condiciones son equivalentes y cómo se transforman uno al otro?
    \item La \textbf{correspondencia entre campo de velocidad y función score}: ¿cuál es la expresión matemática exacta entre velocity y score function?
    \item \textbf{Rectified Flow}: ¿cómo simplifica el proceso de generación las trayectorias rectas y cómo permite el ajuste fino con Reflow hacerlas aún más lineales?
    \item Perspectiva de \textbf{transporte óptimo}: ¿cómo afectan los emparejamientos aleatorios frente a los óptimos al grado de curvatura de las trayectorias?
    \item \textbf{Aplicaciones prácticas}: ¿cómo aplican las tecnologías de Flow Matching los modelos Stable Diffusion 3 (SD3) y Flux?
\end{enumerate}

\begin{knowledgebox}{Repaso de conocimientos previos}
Para comprender esta clase se requieren los siguientes conceptos previos introducidos en la clase anterior (Lecture 9: Flow Matching 1):
\begin{itemize}
    \item \textbf{Mapa de flujo (Flow Map)} $\psi_t$: mapeo determinista desde $t=0$ hasta $t$, con $\psi_0(\mathbf{x}) = \mathbf{x}$
    \item \textbf{Campo vectorial (Vector Field)} $u_t(\mathbf{x})$: derivada temporal del mapa de flujo, $\frac{d\psi_t}{dt} = u_t(\psi_t(\mathbf{x}))$
    \item \textbf{Trayectoria de probabilidad (Probability Path)} $p_t$: distribución de probabilidad que varía con el tiempo $t$, satisfaciendo $p_0 = p_{\text{ref}}$ y $p_1 = p_{\text{data}}$
    \item \textbf{Ecuación de Fokker--Planck}: ecuación que vincula el campo vectorial con la trayectoria de probabilidad
    \item \textbf{Flow Matching condicional (Conditional Flow Matching)}: descomposición del campo vectorial marginal en una expectativa sobre campos vectoriales condicionales
\end{itemize}
\end{knowledgebox}

\subsection{Convenciones de notación y cambio del dominio temporal}

En el contexto de Flow Matching, el dominio temporal difiere del de los modelos de difusión:

\begin{itemize}
    \item \textbf{Modelos de difusión}: el tiempo va desde $T$ (extremo ruidoso) hasta $0$ (extremo de datos); la generación es un ``proceso inverso''
    \item \textbf{Modelos de flujo}: el tiempo va desde $0$ (distribución de referencia) hasta $1$ (distribución de datos); la generación es un ``proceso directo''
\end{itemize}

Específicamente:
\begin{itemize}
    \item $\mathbf{x}_0 \sim p_0 = \mathcal{N}(\mathbf{0}, \mathbf{I})$: muestra proveniente de la distribución de referencia (gausiana estándar)
    \item $\mathbf{x}_1 \sim p_1 = p_{\text{data}}$: muestra proveniente de la distribución de datos
\end{itemize}

\begin{warningbox}{Símbolos propensos a confusión}
Observe que en Flow Matching, $\mathbf{x}_0$ es el \textbf{ruido} y $\mathbf{x}_1$ son los \textbf{datos}, lo cual es exactamente opuesto a la convención habitual en DDPM donde $\mathbf{x}_0$ son datos y $\mathbf{x}_T$ es ruido. Además, las funciones $\mu_t, \sigma_t$ de la Lecture 9 son esencialmente idénticas a los $\alpha_t, \sigma_t$ de los modelos de difusión variacional; la única diferencia radica en el cambio del dominio temporal, que invierte sus posiciones relativas.
\end{warningbox}

\subsection{Forma general del mapa de flujo condicional}

En la clase anterior se introdujo la forma lineal gaussiana del mapa de flujo condicional (Conditional Flow Map):
$$\psi_t(\mathbf{x}_0 \mid \mathbf{x}_1) = \sigma_t \cdot \mathbf{x}_0 + \alpha_t \cdot \mathbf{x}_1$$

Donde:
\begin{itemize}
    \item $\sigma_t, \alpha_t$ son funciones escalares del tiempo $t$
    \item Condiciones de frontera: $\sigma_0 = 1, \alpha_0 = 0$ (en $t=0$ se obtiene puro ruido $\mathbf{x}_0$); $\sigma_1 = 0, \alpha_1 = 1$ (en $t=1$ se obtienen datos $\mathbf{x}_1$)
\end{itemize}

La trayectoria de probabilidad condicional correspondiente es una distribución gaussiana:
$$p_t(\mathbf{x} \mid \mathbf{x}_1) = \mathcal{N}(\mathbf{x}; \alpha_t \mathbf{x}_1, \sigma_t^2 \mathbf{I})$$

Tomando la derivada temporal del mapa de flujo condicional se obtiene el campo vectorial condicional:
$$u_t(\mathbf{x} \mid \mathbf{x}_1) = \dot{\sigma}_t \cdot \mathbf{x}_0 + \dot{\alpha}_t \cdot \mathbf{x}_1$$

\subsection{Resumen de la sección}

En esta sección se repasan los supuestos básicos de Flow Matching, incluyendo las convenciones del dominio temporal, la definición de mapa de flujo condicional y la derivación del campo vectorial condicional. El punto clave es que en el modelo de flujo el tiempo va desde $0$ (ruido) hasta $1$ (datos), y el objetivo del entrenamiento es aprender una red neuronal para predecir el campo vectorial (velocidad); esto se asemeja conceptualmente a la predicción de ruido en los modelos de difusión.

\section{Equivalencia matemática entre modelos de difusión y modelos de flujo}

Esta sección es la parte teórica central de esta clase, donde el ponente deriva detalladamente la relación matemática exacta entre los modelos de difusión y los modelos de flujo.

\subsection{Partiendo del modelo de difusión variacional}

Recuerde que en el Modelo de Difusión Variacional (Variational Diffusion Model, VDM), el proceso hacia adelante se define como:
$$q(\mathbf{x}_t \mid \mathbf{x}_0) = \mathcal{N}(\mathbf{x}_t; \alpha_t \mathbf{x}_0, \sigma_t^2 \mathbf{I})$$

En el contexto de Flow Matching, interprete este proceso hacia adelante nuevamente como un mapeo de flujo condicional:
$$\psi_t(\mathbf{x}_0 \mid \mathbf{x}_1) = \sigma_t \cdot \mathbf{x}_0 + \alpha_t \cdot \mathbf{x}_1$$

\begin{importantbox}{Insight central: los modelos de difusión son un caso especial de modelos de flujo}
Cuando el mapeo de flujo condicional se define en la forma gaussiana lineal anterior y las funciones $\alpha_t, \sigma_t$ coinciden exactamente con las funciones de programación del modelo de difusión variacional, el modelo de difusión puede considerarse un\textbf{caso especial} de los modelos de flujo. En otras palabras, los modelos de flujo son una\textbf{generalización} de los modelos de difusión. La única diferencia entre ambos radica en:
\begin{itemize}
    \item \textbf{Modelos de difusión}: primero se define la trayectoria de probabilidad (proceso de adición de ruido hacia adelante), luego se derivan el campo vectorial y las trayectorias
    \item \textbf{Modelos de flujo}: primero se define el mapeo de flujo condicional (trayectoria), luego se derivan el campo vectorial y la trayectoria de probabilidad
\end{itemize}
\end{importantbox}

\subsection{Reaparecimiento de la fórmula de Tweedie}

Aquí, el ponente introduce una herramienta de cálculo importante: la fórmula de Tweedie. Dado un punto de datos intermedio $\mathbf{x}_t = \sigma_t \mathbf{x}_0 + \alpha_t \mathbf{x}_1$, donde $\mathbf{x}_0 \sim \mathcal{N}(\mathbf{0}, \mathbf{I})$ y $\mathbf{x}_1 \sim p_{\text{data}}$, podemos calcular la esperanza a posteriori de $\mathbf{x}_1$:

$$\mathbb{E}[\mathbf{x}_1 \mid \mathbf{x}_t] = \frac{\mathbf{x}_t - \sigma_t \cdot \mathbb{E}[\mathbf{x}_0 \mid \mathbf{x}_t]}{\alpha_t} = \frac{\mathbf{x}_t + \sigma_t^2 \nabla_{\mathbf{x}_t} \log p_t(\mathbf{x}_t)}{\alpha_t}$$

Donde:
\begin{itemize}
    \item $\nabla_{\mathbf{x}_t} \log p_t(\mathbf{x}_t)$ es la función de puntaje (Score Function)
    \item Esto coincide exactamente con la fórmula de Tweedie presentada en la Clase 5, solo que se utilizan las convenciones de notación de Flow Matching
\end{itemize}

\begin{knowledgebox}{Intuición detrás de la fórmula de Tweedie}
La idea central de la fórmula de Tweedie es: dado una observación $\mathbf{x}_t$ contaminada con ruido, la media a posteriori de la señal original $\mathbf{x}_1$ puede calcularse mediante la función de puntaje. La función de puntaje apunta en la dirección donde crece más rápido la densidad de probabilidad del logaritmo, es decir, la dirección de ``mayor concentración de datos''. En un modelo de verosimilitud gaussiana, esto equivale a moverse desde la posición actual hacia la dirección de la media a posteriori.
\end{knowledgebox}

\subsection{Relación entre el campo de velocidad y la función de puntaje}

Esta es una de las derivaciones más importantes de esta clase. Partiendo del campo vectorial condicional:

$$u_t(\mathbf{x} \mid \mathbf{x}_1) = \dot{\sigma}_t \cdot \mathbf{x}_0 + \dot{\alpha}_t \cdot \mathbf{x}_1$$

Expresando $\mathbf{x}_0$ en términos de $\mathbf{x}_t$ y $\mathbf{x}_1$: $\mathbf{x}_0 = \frac{\mathbf{x}_t - \alpha_t \mathbf{x}_1}{\sigma_t}$, y sustituyendo obtenemos:

$$u_t(\mathbf{x}_t \mid \mathbf{x}_1) = \frac{\dot{\sigma}_t}{\sigma_t}(\mathbf{x}_t - \alpha_t \mathbf{x}_1) + \dot{\alpha}_t \mathbf{x}_1$$

$$= \frac{\dot{\sigma}_t}{\sigma_t} \mathbf{x}_t + \left(\dot{\alpha}_t - \frac{\dot{\sigma}_t \alpha_t}{\sigma_t}\right) \mathbf{x}_1$$

Tomando la esperanza sobre $\mathbf{x}_1$ (marginalizando) y utilizando la fórmula de Tweedie, obtenemos el campo vectorial marginal:

\begin{importantbox}{Fórmula equivalente entre el campo de velocidad y la función de puntaje}
$$u_t(\mathbf{x}_t) = \frac{\dot{\sigma}_t}{\sigma_t} \mathbf{x}_t + \left(\dot{\alpha}_t - \frac{\dot{\sigma}_t \alpha_t}{\sigma_t}\right) \mathbb{E}[\mathbf{x}_1 \mid \mathbf{x}_t]$$

Sustituyendo la fórmula de Tweedie en $\mathbb{E}[\mathbf{x}_1 \mid \mathbf{x}_t]$, se puede establecer una relación equivalente entre el\textbf{campo de velocidad} y la\textbf{función de puntaje}:
$$u_t(\mathbf{x}_t) = \frac{\dot{\sigma}_t}{\sigma_t} \mathbf{x}_t + \left(\dot{\alpha}_t - \frac{\dot{\sigma}_t \alpha_t}{\sigma_t}\right) \cdot \frac{\mathbf{x}_t + \sigma_t^2 \nabla_{\mathbf{x}_t} \log p_t(\mathbf{x}_t)}{\alpha_t}$$

Esto significa que: dado cualquier punto de datos intermedio $\mathbf{x}_t$ y un tiempo $t$, basta con conocer la función de puntaje para calcular el campo de velocidad, y viceversa.
\end{importantbox}

\subsection{Relación entre el campo de velocidad y la predicción de ruido}

Utilizando la relación conocida entre la función de puntaje y la predicción de ruido $\nabla_{\mathbf{x}_t} \log p_t(\mathbf{x}_t) = -\frac{\boldsymbol{\epsilon}}{\sigma_t}$, podemos establecer aún más la relación equivalente entre el campo de velocidad y la red de predicción de ruido:

$$u_t(\mathbf{x}_t) = \frac{\dot{\alpha}_t}{\alpha_t} \mathbf{x}_t - \left(\dot{\alpha}_t \frac{\sigma_t}{\alpha_t} - \dot{\sigma}_t\right) \boldsymbol{\epsilon}_\theta(\mathbf{x}_t, t)$$

Donde $\boldsymbol{\epsilon}_\theta$ es la red entrenada de predicción de ruido en el modelo de difusión.

\begin{warningbox}{Predicción de velocidad vs.\ predicción de ruido: no es solo cambiar la cabeza de salida}
Aunque el campo de velocidad y el campo de ruido pueden transformarse exactamente entre sí matemáticamente, elegir predecir una u otra cantidad durante el\textbf{entrenamiento} afecta el comportamiento numérico de los gradientes. Las prácticas recientes indican que entrenar directamente la predicción de velocidad (v-prediction) ofrece mayor estabilidad numérica en ciertos esquemas de programación en comparación con la predicción de ruido ($\epsilon$-prediction), especialmente cuando $t$ se acerca a los extremos. Esto es una de las razones por las que SD3 y Flux optaron por v-prediction.
\end{warningbox}

\subsection{Consistencia con el PF-ODE de flujo de probabilidad}

El ponente señala además que el campo vectorial derivado desde Flow Matching es completamente consistente con el Probability Flow ODE (PF-ODE) derivado desde el marco SDE. Recuerde la forma general del PF-ODE:

$$d\mathbf{x} = \left[f(\mathbf{x}, t) - \frac{1}{2}g(t)^2 \nabla_\mathbf{x} \log p_t(\mathbf{x})\right] dt$$

Donde $f(\mathbf{x}, t)$ y $g(t)$ son el término de deriva y el coeficiente de difusión del SDE, respectivamente. Sus relaciones con $\alpha_t, \sigma_t$ son:
$$f(\mathbf{x}, t) = \frac{\dot{\alpha}_t}{\alpha_t} \mathbf{x}, \quad g(t)^2 = \frac{d\sigma_t^2}{dt} - 2\frac{\dot{\alpha}_t}{\alpha_t}\sigma_t^2$$

Al sustituir estos valores en el PF-ODE, se puede verificar que su resultado es\textbf{completamente idéntico} al campo vectorial derivado mediante Flow Matching.

\begin{importantbox}{Unificación de tres perspectivas equivalentes}
Hasta aquí, hemos establecido una relación de equivalencia completa entre las siguientes tres perspectivas (bajo un camino condicional gaussiano lineal):
\begin{enumerate}
    \item \textbf{Perspectiva del modelo de difusión}: aprender la función de puntaje $\nabla_\mathbf{x} \log p_t(\mathbf{x})$ o el ruido $\boldsymbol{\epsilon}$
    \item \textbf{Perspectiva de Flow Matching}: aprender el campo vectorial/velocidad $u_t(\mathbf{x})$
    \item \textbf{Perspectiva del PF-ODE}: obtener una trayectoria de muestreo determinista mediante la conversión de SDE a ODE
\end{enumerate}
Las tres son transformables exactamente entre sí en términos matemáticos; las diferencias radican en qué cantidad se predice durante el entrenamiento y qué estrategia de resolución se utiliza al momento del muestreo.
\end{importantbox}

\subsection{Resumen de la sección}

Esta sección ha establecido una relación matemática exacta de equivalencia entre los modelos de difusión y los modelos de flujo. La conclusión central es que, bajo un mapeo de flujo condicional gaussiano lineal, el modelo de difusión es un caso especial del modelo de flujo. El campo de velocidad, la función de puntaje y la predicción de ruido pueden intercambiarse libremente. Esta relación de equivalencia no solo tiene significado teórico, sino que implica en la práctica que: un modelo de difusión preentrenado puede ser directamente ``reinterpretado'' como un modelo de flujo para su uso.

\section{Ventajas del modelo de flujo y capacidad de generalización}

Tras establecer la relación de equivalencia, el ponente se dirigió a discutir las \textbf{ventajas únicas} del modelo de flujo frente al modelo de difusión —es decir, en qué aspectos el modelo de flujo realmente ``supera'' al modelo de difusión.

\subsection{Mapa de flujo condicional no lineal}

Los modelos de difusión utilizan esencialmente caminos condicionales de forma gaussiana lineal. En cambio, el marco Flow Matching permite definir \textbf{cualquier} mapa de flujo condicional, sin limitarse a la forma lineal:

\begin{itemize}
    \item Se pueden probar mapas no lineales $\psi_t(\mathbf{x}_0 \mid \mathbf{x}_1)$, por ejemplo añadiendo transformaciones no lineales como coseno o sigmoide
    \item Los experimentos demuestran que, aunque los mapas de flujo no lineales actualmente no superan significativamente la forma lineal, la flexibilidad inherente del marco ofrece espacio para futuras exploraciones
\end{itemize}

\subsection{Distribución de referencia arbitraria}

Esta es una de las generalizaciones más importantes del modelo de flujo. Los modelos de difusión requieren que la distribución de referencia sea necesariamente una distribución gaussiana estándar $\mathcal{N}(\mathbf{0}, \mathbf{I})$, ya que toda la teoría (proceso de difusión hacia adelante) se construye sobre la base del ruido gaussiano. En contraste, el modelo de flujo puede utilizar \textbf{cualquier} distribución de referencia:

\begin{importantbox}{Mapeo de distribución a distribución}
Flow Matching permite aprender un mapeo desde \textbf{cualquier} distribución fuente hacia \textbf{cualquier} distribución objetivo. Esto incluye no solo:
\begin{itemize}
    \item Gaussiana estándar $\to$ distribución de datos (tarea de generación clásica)
\end{itemize}
sino también:
\begin{itemize}
    \item Distribución de datos A $\to$ distribución de datos B (transferencia de estilo, conversión de dominios)
    \item Distribución de baja resolución $\to$ distribución de alta resolución (superresolución)
    \item Conversiones entre distribuciones arbitrariamente complejas
\end{itemize}
Esta flexibilidad es algo que el marco del modelo de difusión no posee.
\end{importantbox}

\subsection{Origen ODE vs.\ SDE}

\begin{knowledgebox}{Dos filosofías de modelado}
Otra diferencia fundamental entre los modelos de difusión y los de flujo radica en el \textbf{origen} del modelado:
\begin{itemize}
    \item \textbf{Modelos de difusión}: parten de una \textbf{ecuación diferencial estocástica} (SDE), definiendo el proceso de difusión hacia adelante como un proceso estocástico que añade ruido gaussiano gradualmente; luego se descubre que una ODE \textbf{determinista} equivalente también proporciona la misma distribución marginal
    \item \textbf{Modelos de flujo}: parten directamente de una \textbf{ecuación diferencial ordinaria} (ODE), definiendo un mapeo determinista desde la distribución fuente hacia la distribución objetivo, sin necesidad de introducir aleatoriedad
\end{itemize}
El origen del modelo de flujo es más sencillo y evita las derivaciones complejas de conversión de SDE a ODE.
\end{knowledgebox}

\subsection{Resumen de la sección}

El modelo de flujo generaliza al modelo de difusión en los siguientes aspectos: (1) el mapa de flujo condicional no se limita a formas lineales; (2) la distribución de referencia no se limita a una gaussiana estándar; (3) parte directamente de una ODE, ofreciendo un concepto más sencillo. Estas generalizaciones proporcionan una base teórica para el diseño de modelos generadores más flexibles.
\section{Rectified Flow y trayectorias rectas}

Esta sección es la parte más importante de esta clase desde el punto de vista de su aplicación; el conferenciante detalló exhaustivamente las ideas de Rectified Flow (flujo corregido), sus métodos de entrenamiento y la técnica para enderezar las trayectorias mediante Reflow.

\subsection{Definición de Rectified Flow}

Rectified Flow es un caso especial de Flow Matching, donde el mapeo del flujo condicional se define como una \textbf{interpolación lineal} entre $\mathbf{x}_0$ y $\mathbf{x}_1$:

$$\psi_t(\mathbf{x}_0 \mid \mathbf{x}_1) = (1 - t)\, \mathbf{x}_0 + t\, \mathbf{x}_1$$

El campo vectorial condicional (velocidad) correspondiente es extremadamente conciso:

$$u_t(\mathbf{x}_t \mid \mathbf{x}_0, \mathbf{x}_1) = \mathbf{x}_1 - \mathbf{x}_0$$

Es decir: la velocidad condicional es constante y su dirección siempre apunta en línea recta desde $\mathbf{x}_0$ hacia $\mathbf{x}_1$.

\begin{importantbox}{Proceso de entrenamiento de Rectified Flow}
El proceso de entrenamiento de Rectified Flow es extremadamente conciso:
\begin{enumerate}
    \item Muestrear del distribución de datos $\mathbf{x}_1 \sim p_{\text{data}}$
    \item Muestrear de la distribución de referencia $\mathbf{x}_0 \sim \mathcal{N}(\mathbf{0}, \mathbf{I})$
    \item Muestrear uniformemente el tiempo $t \sim \mathcal{U}(0, 1)$
    \item Calcular el punto intermedio $\mathbf{x}_t = (1-t)\mathbf{x}_0 + t\mathbf{x}_1$
    \item Calcular la velocidad objetivo $v^* = \mathbf{x}_1 - \mathbf{x}_0$
    \item Minimizar $\mathcal{L} = \| v_\theta(\mathbf{x}_t, t) - v^* \|^2$
\end{enumerate}
\end{importantbox}

\subsection{Trayectorias condicionales vs.\ trayectorias marginales}

Este es el punto clave para comprender las limitaciones de Rectified Flow. Aunque cada par $(\mathbf{x}_0, \mathbf{x}_1)$ tiene una \textbf{trayectoria condicional} estrictamente recta, las trayectorias generadas por el \textbf{campo vectorial marginal} que produce el modelo después del entrenamiento generalmente \textbf{no son} rectas.

\begin{warningbox}{Las líneas rectas condicionales no equivalen a líneas rectas marginales}
Consideremos un ejemplo bidimensional: la distribución de origen tiene dos grupos gaussianos, y la distribución objetivo también tiene dos grupos gaussianos. La trayectoria condicional entre cada par $(\mathbf{x}_0, \mathbf{x}_1)$ es efectivamente una línea recta; sin embargo, dado que $\mathbf{x}_0$ y $\mathbf{x}_1$ se emparejan de manera \textbf{aleatoria e independiente}, en ciertos puntos intermedios $\mathbf{x}_t$, la velocidad condicional podría apuntar hacia direcciones muy diferentes (hacia arriba o hacia abajo). Después del entrenamiento, la red neuronal devuelve el \textbf{promedio} (esperanza) de estas velocidades condicionales; como resultado, puede que ni apunte hacia arriba ni hacia abajo, sino hacia una dirección promedio entre ambas. Por lo tanto, la trayectoria marginal se curva.
\end{warningbox}

**Consecuencias de las trayectorias curvas**: si la trayectoria marginal es curva, se requieren suficientes pasos temporales al resolver la EDO para rastrear con precisión la trayectoria. Cuantos menos pasos, mayor será el error de discretización y peor la calidad generada. En el caso ideal, si todas las trayectorias marginales fueran rectas, bastaría con un \textbf{único paso} para ir directamente de $\mathbf{x}_0$ a $\mathbf{x}_1$.

\subsection{Reflow: enderezar las trayectorias mediante ajuste fino}

Para hacer que las trayectorias marginales se acerquen más a líneas rectas, Liu et al.\ (2023) propusieron la técnica de \textbf{Reflow} (reflujo):

\begin{enumerate}
    \item \textbf{Fase 1}: Entrenar normalmente un modelo de Rectified Flow (1-Rectified Flow)
    \item \textbf{Generar datos emparejados}: Usando el modelo entrenado, partir de $\mathbf{x}_0 \sim \mathcal{N}(\mathbf{0}, \mathbf{I})$ y generar $\mathbf{x}_1$ resolviendo completamente la EDO; registrar los pares $(\mathbf{x}_0, \mathbf{x}_1)$
    \item \textbf{Fase 2}: Volver a entrenar el modelo con estos datos emparejados (2-Rectified Flow)
    \item Se puede continuar iterando: generar nuevos pares con 2-Rectified Flow y entrenar 3-Rectified Flow, y así sucesivamente
\end{enumerate}

\begin{importantbox}{¿Por qué funciona Reflow?}
La intuición central de por qué Reflow logra enderezar las trayectorias es:
\begin{itemize}
    \item En la fase 1, $\mathbf{x}_0$ y $\mathbf{x}_1$ son emparejamientos aleatorios muestreados de manera \textbf{independente}, sin relación correspondiente
    \item Después del entrenamiento en la fase 1, el modelo ya ha aprendido un mapeo (curvo) de $\mathbf{x}_0$ a $\mathbf{x}_1$
    \item En la fase de Reflow, los pares $(\mathbf{x}_0, \mathbf{x}_1)$ ya no son aleatorios; provienen de un \textbf{mapeo determinista} definido por el modelo. Ya existe una relación correspondiente razonable entre estos pares
    \item Volver a entrenar con estos ``pares significativos'' reduce drásticamente las intersecciones y conflictos de las trayectorias, haciendo que la trayectoria marginal se enderece naturalmente
\end{itemize}
\end{importantbox}

Los experimentos demuestran:
\begin{itemize}
    \item \textbf{1-Rectified Flow}: la trayectoria marginal muestra una curvatura notable
    \item \textbf{2-Rectified Flow}: la trayectoria se endereza significativamente
    \item \textbf{3-Rectified Flow}: la trayectoria es casi una línea recta
\end{itemize}

\subsection{Trayectorias rectas y generación con pocos pasos}

La ventaja directa de las trayectorias rectas es que permite realizar la generación con un número \textbf{muy reducido} de pasos de resolución de EDO:

\begin{itemize}
    \item Si la trayectoria es una línea recta perfecta, basta predecir una vez la velocidad $v_\theta(\mathbf{x}_0, 0)$ en el punto inicial $\mathbf{x}_0$ y avanzar en esa dirección hasta el final; esto constituye la \textbf{generación de un solo paso}
    \item Incluso si la trayectoria se acerca a ser recta pero no es perfecta, una resolución de Euler con 2--4 pasos obtiene resultados de alta calidad
    \item Comparación: los modelos de difusión estándar suelen requerir 20--100 pasos para obtener una calidad equivalente
\end{itemize}

\begin{knowledgebox}{De GAN a difusión y luego a modelos de flujo de un solo paso}
El conferenciante señaló un ciclo interesante en el desarrollo de los modelos generativos:
\begin{itemize}
    \item \textbf{Era de GAN/VAE}: generación de un solo paso, rápida pero limitada en calidad/diversidad
    \item \textbf{Era de los modelos de difusión}: iteración multietapa, con una mejora sustancial en la calidad a costa de la velocidad
    \item \textbf{Rectified Flow + Reflow}: mediante el ajuste fino se regresa a la generación de un solo paso o con pocos pasos, manteniendo al mismo tiempo una alta calidad
\end{itemize}
Se trata de una vía de desarrollo que asciende en forma de ``hélice''.
\end{knowledgebox}

\subsection{Resumen de la sección}

Rectified Flow define el mapeo del flujo mediante interpolación lineal, lo que hace que el entrenamiento sea extremadamente conciso. Sin embargo, los emparejamientos aleatorios provocan que la trayectoria marginal se curve, requiriendo múltiples pasos de resolución de EDO. La técnica Reflow logra ``enderezar'' eficazmente la trayectoria marginal utilizando datos generados por el propio modelo para el ajuste fino, lo que posibilita una generación con pocos pasos e incluso con un solo paso.
\section{Perspectiva del transporte óptimo}

La causa fundamental de la curvatura de las trayectorias radica en la \textbf{forma de emparejamiento} de $\mathbf{x}_0$ y $\mathbf{x}_1$. En esta sección se aborda este problema desde la perspectiva del transporte óptimo (Optimal Transport).

\subsection{Emparejamiento aleatorio vs.\ emparejamiento óptimo}

En el entrenamiento estándar de Flow Matching, $\mathbf{x}_0$ y $\mathbf{x}_1$ se muestrean de forma \textbf{independiente}:
$$\mathbf{x}_0 \sim p_0, \quad \mathbf{x}_1 \sim p_1, \quad \text{ambos independientes}$$

Esto significa que una muestra ruidosa $\mathbf{x}_0$ podría ser ``asignada'' a cualquier muestra de datos $\mathbf{x}_1$, sin preferencia alguna por la estructura espacial.

\begin{importantbox}{Idea fundamental del transporte óptimo (Optimal Transport)}
El problema del transporte óptimo se centra en: ¿cómo encontrar el plan de transporte más ``económico'' desde la distribución $p_0$ hasta la distribución $p_1$? Específicamente, para un costo $L_2$ (también conocido como distancia Wasserstein-2), el plan de transporte óptimo minimiza:
$$W_2^2(p_0, p_1) = \inf_{\pi \in \Pi(p_0, p_1)} \int \|\mathbf{x}_0 - \mathbf{x}_1\|^2 \, d\pi(\mathbf{x}_0, \mathbf{x}_1)$$
donde $\Pi(p_0, p_1)$ es el conjunto de todos los acoplamientos (distribuciones conjuntas), cuyas márgenes son $p_0$ y $p_1$, respectivamente.

\medskip
\textbf{Intuición}: El transporte óptimo tiende a emparejar puntos ``cercanos'', evitando desplazamientos de larga distancia, lo que minimiza así el costo total del traslado.
\end{importantbox}

\subsection{Influencia del emparejamiento OT en la rectitud de las trayectorias}

Si utilizamos un plan de transporte óptimo para emparejar $\mathbf{x}_0$ y $\mathbf{x}_1$, en lugar de un emparejamiento aleatorio, entonces:

\begin{itemize}
    \item El emparejamiento es más ``racional'': los puntos ruidosos cercanos tienden a mapearse hacia puntos de datos cercanos.
    \item Hay menos cruces entre las trayectorias condicionales de distintos emparejamientos.
    \item Las colisiones en los campos vectoriales marginales son menores, lo que hace que las trayectorias marginales sean naturalmente más rectas.
\end{itemize}

\begin{warningbox}{Costo computacional del emparejamiento OT}
Aunque el emparejamiento OT es teóricamente muy atractivo, la complejidad de calcular exactamente un plan de transporte óptimo es muy alta. Para $n$ muestras, la complejidad de los algoritmos clásicos es $O(n^3)$, lo cual no es práctico en entrenamientos a gran escala. Por lo tanto, en la práctica se suelen emplear métodos aproximados, como calcular OT dentro de un mini-lote (Mini-batch OT), o utilizar Reflow como alternativa; Reflow puede considerarse una optimización de emparejamiento OT ``implícita''.
\end{warningbox}

\subsection{Relación entre Reflow y OT}

Existe una conexión profunda entre Reflow y el transporte óptimo:

\begin{itemize}
    \item Reflow optimiza implícitamente el plan de emparejamiento $\mathbf{x}_0 \to \mathbf{x}_1$ mediante ajuste fino iterativo.
    \item Cada ronda de Reflow acerca el emparejamiento al plan OT, ya que el mapeo aprendido por el modelo tiende a seguir la ``ruta más corta''.
    \item Teóricamente, el límite de infinitas iteraciones de Reflow es el mapeo OT (bajo ciertas condiciones de regularidad).
\end{itemize}

\begin{knowledgebox}{Curvatura de las trayectorias vista desde el costo de transporte}
La esencia de la curvatura de las trayectorias radica en un \textbf{costo de transporte excesivo}. Cuando se emparejan puntos distantes, la línea que los conecta atraviesa regiones intermedias, cruzándose con otras líneas rectas de distintos emparejamientos. Estos cruces generan inconsistencias en el campo vectorial marginal (múltiples direcciones de velocidad en el mismo punto); la red solo puede aprender la velocidad esperada (promedio), lo que resulta en trayectorias curvas. El emparejamiento OT reduce los cruces al minimizar el costo de transporte, mientras que Reflow logra un efecto similar mediante una optimización implícita iterativa.
\end{knowledgebox}

\subsection{Resumen de la sección}

El transporte óptimo proporciona un marco teórico para comprender el problema de la curvatura de las trayectorias. El emparejamiento aleatorio conduce a un alto costo de transporte y cruces en las trayectorias; el emparejamiento OT puede aliviar este problema, aunque su costo computacional es elevado. Reflow ofrece una alternativa práctica que optimiza implícitamente el emparejamiento mediante ajuste fino iterativo.

````
\section{Aplicación práctica: Stable Diffusion 3 y Flux}

En esta sección, el ponente ofrece una breve introducción sobre la aplicación de la tecnología Flow Matching en los modelos actuales más avanzados para generación de imágenes y video.

\subsection{Fundamentos de flujo de modelos para SD3 y Flux}

El ponente señala explícitamente que muchos de los modelos generativos de vanguardia actuales han \textbf{cambiado de} modelos de difusión a modelos de flujo:

\begin{itemize}
    \item \textbf{Stable Diffusion 3} (SD3): Aunque el nombre incluye ``diffusion'', en realidad utiliza entrenamiento basado en Flow Matching; la red predice la velocidad (velocity) en lugar del ruido.
    \item \textbf{Flux}: Modelo desarrollado por Black Forest Labs, que también se basa en la tecnología Rectified Flow.
\end{itemize}

\begin{importantbox}{Elecciones de diseño clave para SD3}
Las decisiones técnicas centrales de Stable Diffusion 3 incluyen:
\begin{itemize}
    \item \textbf{Entrenamiento con Rectified Flow}: Uso de mapeo de flujo condicional mediante interpolación lineal; la red genera una velocidad $v_\theta(\mathbf{x}_t, t)$.
    \item \textbf{Arquitectura MM-DiT}: Adopción de Multimodal Diffusion Transformer como red base.
    \item \textbf{Objetivo de trayectoria recta}: Búsqueda de trayectorias de generación más directas mediante técnicas como Reflow, reduciendo así el número de pasos de muestreo.
    \item \textbf{Predicción de velocidad (v-prediction)}: En lugar de predicción de ruido ($\epsilon$-prediction), lo cual ofrece una mejor estabilidad numérica en ciertos niveles de ruido.
\end{itemize}
\end{importantbox}

\subsection{Modelo Flux}

El ponente menciona que, desde la primera sesión del curso, se solicitó a los estudiantes generar imágenes utilizando el modelo Flux. Una variante clave de Flux es \textbf{Flux.1-schnell}, que emplea la tecnología Rectified Flow junto con Reflow para lograr generación de imágenes de alta calidad en un solo paso o con pocos pasos.

\begin{knowledgebox}{Modelos de consistencia vs.\ Rectified Flow: dos rutas tecnológicas para acelerar la generación}
En cuanto a la aceleración del muestreo en modelos de difusión/flujo, actualmente existen dos rutas competitivas principales:
\begin{itemize}
    \item \textbf{Consistency Model (Modelo de consistencia)}: Aprendizaje de un ``mapeo rápido'' sobre las trayectorias ODE mediante Consistency Distillation (CD) o Consistency Training (CT), permitiendo saltar directamente desde cualquier punto intermedio al final. Esto se discutió en la Lecture 8/9.
    \item \textbf{Rectified Flow + Reflow}: Reducción del número de pasos necesarios al hacer que las trayectorias mismas sean más rectas. Este es el tema central de esta sesión.
\end{itemize}

Los resultados experimentales comparativos muestran:
\begin{itemize}
    \item El Consistency Model, al generar en 1--2 pasos, alcanza un FID aproximado de 3--4 en ImageNet 64$\times$64.
    \item El 2-Rectified Flow también logra un FID similar en generación de 1--2 pasos.
\end{itemize}
Ambos métodos tienen sus propias ventajas y desventajas, y siguen siendo áreas de investigación activa.
\end{knowledgebox}

\subsection{Resumen de la sección}

El Flow Matching, especialmente las variantes de Rectified Flow, se ha convertido en la tecnología fundamental para los modelos generadores de imágenes más avanzados del momento (SD3, Flux). Su ventaja central radica en un entrenamiento sencillo, trayectorias controlables y la posibilidad de generar con pocos pasos.


\section{Comparación de los flujos completos de entrenamiento e inferencia}

Para ayudar a los lectores a establecer una comprensión integral, en esta sección se contrasta mediante tablas las diferencias entre los modelos difusivos y los modelos de flujo en cada etapa del entrenamiento y la inferencia.

\subsection{Comparación de la fase de entrenamiento}

\begin{table}[H]
\centering
\caption{Modelos difusivos vs.\ modelos de flujo: comparación de la fase de entrenamiento}
\begin{tabular}{@{} p{3.8cm} p{5cm} p{5cm} @{}}
\toprule
\textbf{Etapa} & \textbf{Modelo difusivo} & \textbf{Modelo de flujo (Rectified Flow)} \\
\midrule
Dominio temporal & $t \in \{0, \ldots, T\}$ o ${[}0, T{]}$ & $t \in {[}0, 1{]}$ \\
Proceso hacia adelante & $\mathbf{x}_t = \sqrt{\bar\alpha_t}\,\mathbf{x}_0 + \sqrt{1-\bar\alpha_t}\,\boldsymbol{\epsilon}$ & $\mathbf{x}_t = (1-t)\mathbf{x}_0 + t\mathbf{x}_1$ \\
Objetivo de predicción & Ruido $\boldsymbol{\epsilon}$ o gradiente logarítmico $\nabla\log p_t$ & Velocidad $v = \mathbf{x}_1 - \mathbf{x}_0$ \\
Función de pérdida de entrenamiento & $\|\boldsymbol{\epsilon}_\theta(\mathbf{x}_t,t) - \boldsymbol{\epsilon}\|^2$ & $\|v_\theta(\mathbf{x}_t,t) - (\mathbf{x}_1-\mathbf{x}_0)\|^2$ \\
Requerimientos de datos & Solo se necesita $\mathbf{x}_0 \sim p_{\text{data}}$ & Se requiere $\mathbf{x}_0 \sim p_0$ y $\mathbf{x}_1 \sim p_1$ \\
\bottomrule
\end{tabular}
\end{table}

\subsection{Comparación de la fase de inferencia}

\begin{table}[H]
\centering
\caption{Modelos difusivos vs.\ modelos de flujo: comparación de la fase de inferencia}
\begin{tabular}{@{} p{3.8cm} p{5cm} p{5cm} @{}}
\toprule
\textbf{Etapa} & \textbf{Modelo difusivo} & \textbf{Modelo de flujo} \\
\midrule
Punto inicial & $\mathbf{x}_T \sim \mathcal{N}(\mathbf{0},\mathbf{I})$ & $\mathbf{x}_0 \sim \mathcal{N}(\mathbf{0},\mathbf{I})$ \\
Método de resolución & Ecuación diferencial estocástica (SDE) o ODE con función de trayectoria (PF-ODE) & ODE \\
Actualización por paso & $\mathbf{x}_{t-1} = \mathbf{x}_t + f(\mathbf{x}_t,t)\Delta t$ & $\mathbf{x}_{t+\Delta t} = \mathbf{x}_t + v_\theta(\mathbf{x}_t,t)\Delta t$ \\
Número típico de pasos & 20--100 pasos & 1--10 pasos (después de Reflow) \\
Resolutores avanzados & DPM-Solver, entre otros & Euler / Midpoint son suficientes \\
\bottomrule
\end{tabular}
\end{table}

\subsection{Resumen de la sección}

Los modelos de flujo ofrecen un proceso más conciso tanto en la fase de entrenamiento como en la de inferencia. Es especialmente digno de notar que, después del paso Reflow, la fase de inferencia del Rectified Flow puede completarse con muy pocos pasos; esto representa su mayor ventaja para el despliegue práctico.
\section{Visión general del sistema de conocimientos del curso}

En la conclusión del curso, el ponente repasa la evolución del conocimiento acumulado hasta el momento.

\subsection{Trayectoria de evolución tecnológica}

Desde el inicio del curso hasta la décima clase, el ponente ha construido la siguiente trayectoria tecnológica completa:

\begin{enumerate}
    \item \textbf{VAE/GAN} (Lecture 1--2): modelos generativos de un solo paso, limitados en calidad y diversidad.
    \item \textbf{DDPM} (Lecture 3--4): introducción del desruido iterativo multietapa, lo que eleva drásticamente la calidad generativa.
    \item \textbf{DDIM} (Lecture 5--6): variante de muestreo determinista de DDPM, compatible con saltos en pasos.
    \item \textbf{SDE/ODE} (Lecture 7): marco de tiempo continuo, perspectiva unificada.
    \item \textbf{DPM-Solver} (Lecture 8): aceleración del muestreo aprovechando la estructura ODE.
    \item \textbf{Consistency Model} (Lecture 9): aprendizaje de mapas rápidos para lograr generación con pocos pasos.
    \item \textbf{Flow Matching} (Lecture 9--10): marco más general; Rectified Flow + Reflow logra trayectorias rectas y generación en un paso.
\end{enumerate}

\begin{warningbox}{Flow Matching no es una ``alternativa'' a los modelos de difusión}
Aunque esta clase destaca las ventajas de los modelos de flujo, es necesario observar lo siguiente:
\begin{itemize}
    \item Los modelos de difusión y los modelos de flujo son matemáticamente equivalentes (bajo trayectorias condicionales gaussianas lineales).
    \item Un modelo de difusión ya entrenado puede convertirse directamente en un modelo de flujo para su uso.
    \item La elección del marco depende principalmente de consideraciones prácticas: estabilidad del entrenamiento, eficiencia del muestreo y complejidad del código.
    \item Las ``mejores prácticas'' actuales consisten en utilizar el marco Flow Matching + entrenamiento con Rectified Flow + ajuste fino (fine-tuning) con Reflow.
\end{itemize}
\end{warningbox}

\subsection{Próximamente: contenido de las siguientes clases}

El ponente anuncia el contenido de las tres próximas clases:
\begin{itemize}
    \item \textbf{Próxima semana}: Inference-time Guidance (guía en tiempo de inferencia), cómo inyectar preferencias del usuario sin volver a entrenar.
    \item \textbf{Posteriormente}: Discrete Diffusion Models (modelos de difusión para datos discretos), extensión de las ideas de difusión/flujo a datos no continuos (por ejemplo, texto).
\end{itemize}

\subsection{Resumen de la sección}

Desde GAN hasta los modelos de difusión y luego los modelos de flujo, el campo de los modelos generativos ha experimentado una evolución desde un paso, pasando por múltiples pasos, hasta pocos pasos. Flow Matching ofrece un marco unificado y más flexible que permite comprender los logros existentes de los modelos de difusión al mismo tiempo que abre espacio para mejoras futuras.
\section{Síntesis y ampliación}

\subsection{Repaso de los puntos clave}

El aporte central de esta clase (Flow Matching 2) radica en establecer un \textbf{puente matemático completo} entre los modelos de difusión y los modelos de flujo, y en demostrar cómo utilizar dicho puente para diseñar modelos generadores más eficientes. Las conclusiones principales son:

\begin{enumerate}
    \item \textbf{Equivalencia matemática}: Bajo trayectorias condicionales gaussianas lineales, los modelos de difusión son un caso particular de los modelos de flujo. El campo de velocidades $u_t$, la función del gradiente logarítmico $\nabla \log p_t$ y la predicción de ruido $\boldsymbol{\epsilon}_\theta$ pueden convertirse uno en otro con precisión exacta.

    \item \textbf{Unificación de tres perspectivas}: SDE (difusión), PF-ODE (flujo de probabilidad) y Flow ODE (ajuste de flujo) son matemáticamente idénticas.

    \item \textbf{Rectified Flow}: La interpolación lineal define trayectorias condicionales; el objetivo del entrenamiento es la predicción de velocidad, con una formulación extremadamente concisa.

    \item \textbf{Problema de curvatura de las trayectorias}: El emparejamiento aleatorio provoca que las trayectorias condicionales se crucen y que las trayectorias marginales se curve, lo cual requiere resolver mediante ODEs de múltiples pasos.

    \item \textbf{Fine-tuning Reflow}: Mediante la reentrenamiento con datos de emparejamiento generados por el propio modelo, es posible alinear eficazmente las trayectorias y lograr generación con pocos pasos.

    \item \textbf{Conexión con Optimal Transport}: El emparejamiento OT reduce los cruces de trayectoria; Reflow aproxima implícitamente la transformación OT.

    \item \textbf{Impacto práctico}: Modelos actuales líderes como SD3 y Flux ya adoptan Flow Matching / Rectified Flow como tecnología central.
\end{enumerate}

\subsection{Lecturas adicionales}

\begin{itemize}
    \item Lipman, Y., et al. \textit{Flow Matching for Generative Modeling}. ICLR 2023. {\small El artículo fundacional de Flow Matching}
    \item Liu, X., Gong, C., \& Liu, Q. \textit{Flow Straight and Fast: Learning to Generate and Transfer Data with Rectified Flow}. ICLR 2023. {\small Los artículos originales sobre Rectified Flow y Reflow}
    \item Esser, P., et al. \textit{Scaling Rectified Flow Transformers for High-Resolution Image Synthesis}. ICML 2024. {\small Informe técnico de SD3, que extiende Rectified Flow a la generación de imágenes a gran escala}
    \item Tong, A., et al. \textit{Improving and Generalizing Flow-Based Generative Models with Minibatch Optimal Transport}. TMLR 2024. {\small Aplicación de Mini-batch OT en Flow Matching}
    \item Albergo, M. S. \& Vanden-Eijnden, E. \textit{Building Normalizing Flows with Stochastic Interpolants}. ICLR 2023. {\small Perspectiva desde Stochastic Interpolant, estrechamente relacionada con Flow Matching}
    \item Song, Y., et al. \textit{Consistency Models}. ICML 2023. {\small Alternativa competitiva para acelerar el muestreo, comparable a Rectified Flow}
\end{itemize}

%% --- Fin del contenido --- %%

\end{document}
